Define a linear congruential generator to produce pseudorandom numbers. In terms
of quality, recall four issue that should be considered.

A linear congruential generator mimics a uniform distribution on \( [0, 1] \) via a deterministic recursion with the
following structure:
Seed \( x_0 \)
Recursion \( x_{i+1} = (a x_{i} + c) \text{mod} m \) 
Output \( u_{i+1} = \frac{x_{i+1}}{m}\)

This generates random \( U_i \sim \text{unif}[0, 1] \), which are i.i.d.

1. Period length
On finite & deterministic computers the generator eventually repeats the same numbers
Long periods are generally preferred
In the special case with computations modulo \( m \), the period length is at most \( m − 1 \)
2. Reproducibility & portability
It is often important to validate and replicate results. This requires rerunning an algorithm in exactly the same
way, potentially not on the same platform.
3. Speed
4. Randomness

Using the notation of the lecture, for c = 0 state a condition for a full period.

Sufficient criteria for a full period are the following:
1. If \( c \neq 0 \) and if the following three conditions hold, a full period is guaranteed:
Largest common devisor (LcD) of \( c \) and \( m \) is 1.
Every prime number that divides \( m \)m divides \( a - 1 \)
\( a − 1 \) is divisible by \( 4 \) if \( m \) is.
2.
If \( c = 0 \) and if the following three conditions hold, a full period is guaranteed:
1. \( m \) is prime.
2. \( a_{m-1} - 1 \equiv 0 \mod m \)
3. \( a_{j} - 1 \equiv k \mod m \), \( k \neq 0 \), for \( j = 1, 2, \dots, m − 2. \)
If \( a \) satisfies the last two conditions is equivalent to \( a \) is a primitive root of \( m \).

On computers, numbers of different types are typically stored using a fixed number
of bits. Naive implementations of linear congruential generators might fail due to
overflow. Describe two strategies for bounding the size of numbers in intermediate
computations.

Assume there exist natural numbers \( \alpha, a_1, a_2 \) such that
\( a = 2^\alpha a_1 + a_2 \),
\( a_1 , a_2 < 2^\alpha\),
In this case, the recursion of the generator can be represented by numbers that admit smaller upper bounds:
\( ax_i mod m = (a_1 (2^{\alpha} x_i \mod m) + a_2 x_i \mod m) \mod m \)

E.g. assuming \( m = 2^{\beta} - 1 \) then also \( x_i \leq 2^{\beta} - 2 \) and \( (2^{\alpha} x_i \mod m) \leq 2^{\alpha + \beta} - 2^{\alpha + 1} \mod m \),
last multiplying by \( a_1 \) bounds this term strictly by  \( 2^{2 \alpha + \beta} - 2^{2 \alpha + 1} \mod 2^{ \beta } - 1 \)

Assume that
\( a \leq \sqrt{m} \)
Then an algorithm can be specified such that all intermediate calculations result in integers between
\( −(m − 1), \dots, m − 1 \)

Setting \( q = \floor{\frac{m}{a}}, r = m \mod a, m = aq + r \) then the representation is

\( ax_i \mod m = a(x_i \mod q) - \floor{\frac{x_i}{q}}r + (\floor{\frac{x_i}{q}} - \frac{ax_i}{m})m\)

This approach was suggested by P. Bratley, B.L. Fox & L. Schrage (1987).
The specific terms are bounded by
1. \( 0 \leq a (x_i \mod q) < aq ≤ aq + r = m \)
2. \( 0 \leq \floor{\frac{x_i}{q}}r \leq \frac{m-1}{q}r \leq m - 1 \)

The last term does not need to be computed directly. It is \( 0 \), if \( (1) − (2) \geq 0 \), and \( m \), if \( (1) − (2) < 0 \).
This can be seen as follows:
We know that the third summand satisfies \((3) = zm \) for \( z \in \mathbb{Z} \), \( −(m − 1) \leq (1) − (2) \leq m − 1 \), \( 0 < (1) − (2) + (3) \leq m − 1 \) 
(since we are still modulo m)
implying \( z \in \{0, 1\} \) with \( z = 0 \) if \( (1) − (2) \geq 0 \) and \(z = 1\) if \((1) − (2) < 0 \).

The equation itself holds due to

\( ax_i \mod m = ax_i - \floor{\frac{ax_i}{m}}m = (ax_i - \floor{\frac{x_i}{q}}m) + (\floor{\frac{x_i}{q}} - \floor{\frac{ax_i}{m}})m\)
Then the left term can be decomposed as 
\( (ax_i - \floor{\frac{x_i}{q}}m) = ax_i - \floor{\frac{x_i}{q}}(aq + r) = a(x_i - \floor{\frac{x_i}{q}}q) - \floor{\frac{x_i}{q}}r = a(x_i \mod q) - \floor{\frac{x_i}{q}}r \).


Show for \(c = 0\) that linear congruential generators possess a lattice structure.
Discuss the relationship between the distance of the occurring hyperplanes and
the goal of the method.

From 
\(
x_{i+1} = a x_i \mod m
\)

it follows for initial seed \( x_0 \)
\(
x_1 = a x_0 \mod m
\)
\(
x_2 = a x_1 \mod m = a (a x_0 \mod m) \mod m = a^2 x_0 \mod m
\)
Continuing this process, we find that:
\(
x_i = a^i x_0 \mod m
\)
To reveal the lattice structure, we consider the sequence of values \( (x_i, x_{i+1}, x_{i+2}, \ldots, x_{i+k-1}) \) in a k-dimensional space. 
For simplicity, consider the 2-dimensional case \( (x_i, x_{i+1}) \).
\(
(x_i, x_{i+1}) = (a^i x_0 \mod m, a^{i+1} x_0 \mod m)
\)
The key insight is that each point \( (x_i, x_{i+1}) \) lies on the line:
\(
x_{i+1} = a x_i \mod m
\)
If there would be an additional constant involved, the pattern would be less regular since
\(
x_1 = a x_0 + c\mod m
\)
\(
x_2 = a x_1 \mod m = a (a x_0 + c \mod m) \mod m = a^2 x_0 + ac + c\mod m
\)
so 
\(
x_{i} = a^i x_i + a^{i-1}c + a^{i-2}c \dots + c\mod m
\)

Recall the definition of a distribution function and a quantile function. Letting
\( F^{-1} \) be a quantile function and \( U \sim \text{unif}(0, 1) \), derive the 
derive the distribution function of \( F^{-1}(U) \).

Let \( F \) be a cumulative distribution function, and let 
\( F^{-1} \) be its generalized inverse function (using the infimum because CDFs are weakly monotonic and right-continuous)

\( F^{-1}(u) = \inf \{x | F(x) \geq u\} (0<u<1)\)
Claim: If \( U \) is a uniform random variable on \( [0, 1] \) then \( F^{-1}(U) \) has \( F \) as its CDF
\( \mathbb{P}(F^{-1}(U) \leq x) = \mathbb{P}(U \leq F(x)) = F(X) \)
where we used the right-continuity of \( F \) in the first equality.

In case we want to sample a random variable conditional on \( \{a < X \leq b \}\) with \( F(a) < F(b) \),
we use \( V = F(a) + [F(b) - F(a)]U \) then \( F^{-1}(V) \sim \mathcal{L}(X|a < X \leq b)\), since
\( \mathbb{P}(F^{-1}(V) \leq x) = \mathbb{P}(F(a) + [F(b) - F(a)]U \leq F(x))\) 
\( \mathbb{P}(U \leq \frac{F(x) - F(a)}{F(b) - F(a)} = \frac{F(x) - F(a)}{F(b) - F(a)} \).

Illustrate the inverse transform method for the exponential distribution

The exponential distribution is given by \( F(x) = 1 - e^{-x/\theta}, x \geq 0 \)
then \( F^{-1}(u) = - \theta\text{log}(1 - u) \) and \( 1 - u \sim u \) so
\( F^{-1}(u) = - \theta\text{log}(u) \).


Give numerical methods to compute the result of the inverse transform:

If \( F^{-1} \) is analytically tractable (which is more than often not the case)
then it can used directly. 
Otherwise for discrete random variables, one may use a tableau of the 
probability density function e.g. \( \mathbb{P}(X = c_1) = p_{1}, \mathbb{P}(X = c_2) = p_{2} \dots \),
sample a \( U \sim \text{unif}(0, 1) \) and find the closest \( q_i = \sum_{j=1}^{i} p_j \) s.t. \( q_{i-1} U \leq q_i \) and return \( c_i \).

For continuous ones, Newton's method is applicable via fixed point iteration \( F(x) - u = 0 \) to determine a suitable \( x \), e.g. 
\( x_{n+1} = x_n - \frac{F(x_n) - u}{f(x_n)} \).

State the algorithm of the acceptance-rejection-method and prove that it works.

The principle goes as follows:
1. Generate samples from convenient distribution.
2. Reject random subset.
3. Then the accepted samples are distributed according to the target distribution.
To be precise, let \( X ~ f \) be the random variable to be sampled and \( g \) some other density, with \( f(x) \leq cg(x), c > 0 \)

1. Generate \( X \) from \( g \)
2. Generate \( U \sim \text{unif}(0, 1) \)
3. if \( U \leq \frac{f(X)}{cg(X)} \) return \( X \) else goto step \( 1 \).

If let \( Y \) be the output of this algorithm and \( k = \frac{f(X)}{cg(X)} \), then this works because
\( \mathbb{P}(Y \in A) = \mathbb{P}(X \in A|U \leq k) = \frac{\mathbb{P}(X \in A, U \leq k)}{P(U \leq k)} \)
\( = \mathbb{E}[\mathbb{1}_{X \in A}\mathbb{P}(U \leq k|X)](\mathbb{E}[\mathbb{P}(U \leq k|U)])^{-1} = \int_A k g(x) dx (\int_\Omega k g(x) dx = \int_A f(x) dx \)

This can potentially be optimized by using a simpler function than \(\frac{f(X)}{cg(X)} \) for comparison, this is also referred to as the squeeze method.


Explain the acceptance-rejection-algorithm "normal from double exponential".

Double exponential density aka Laplace with \( \mu = 0, b = 1 \) on \( (-\infty, \infty) \), is \( g(x) = \frac{1}{2} e^{-|x|} \)
For \( U_1 \sim \text{unif}(0, 1) \), we have \( -\text{log}(U_1) \sim \text{exp}(1) \) by the inverse transform method,
randomizing the sign with Bernoulli distribution \( p = \frac{1}{2} \), yields a double exponential distribution.
For \( \mathcal{N}(0, 1) \) the density is \( f(x) = \frac{1}{\sqrt{2\pi}}e^{-\frac{1}{2}x^2} \), and with the ratio 
we can find the bounding constant
\( \frac{f(x)}{g(x)} = \sqrt{\frac{2}{\pi}}\exp(-\frac{1}{2}x^2 + |x|) \leq \sqrt{\frac{2e}{\pi}} = c \)
Further \( \frac{f(x)}{cg(x)} = \exp(-\frac{1}{2}x^2 + |x| - \frac{1}{2}) = \exp(-\frac{1}{2}(|x| - 1)^{2} \)

This allows the following simplification of the acceptance-rejection-method
1. Generate \( U_1, U_2, U_3 ~ \text{unif}[0, 1] \)
2. \( X \leftarrow - \text{log}(U_1) \)
3. if \( U_2 > \exp(-\frac{1}{2}(X - 1)^2 \) go to step 1.
4. if \( U_3 \leq \frac{1}{2} \), set \( Y \leftarrow -X \) and return \( Y \)



Explain the Box-Muller-method.
The Box-Muller-method solves the following task, 
generate two independent standard normal random variables \( Z \sim N(0, I_2)\).
note that \( R = Z^2_1 + Z^2_2 \sim \chi_2^{2} = \exp(\frac{1}{2}) \)
and \( \mathcal{(Z_1, Z_2)^{T}|R} = \text{unif}(O) \), where \( O \) denotes the circle of radius \( \sqrt{R} \)

1. Generate \( U_1, U_2 \sim \text{unif}(0, 1) \) independently
2. \( R \leftarrow -2\text{log}(U_1) \)
3. \( V \leftarrow 2\pi U_2 \)
4. \( Z_1 \leftarrow \sqrt{R}\cos(V), Z_2 \leftarrow \sqrt{R}\sin(V)\)
5. return \((Z_1, Z_2)\)

Since \( \sin, \cos \) can be computationally expensive to evaluate, there is an optimization via 
the Marsaglia-Bray-Algorithm.

Explain the Marsaglia-Bray-algorithm.

1. Generate \( U_1, U_2 \sim \text{unif}([-1, 1]) \)
2. \( X \leftarrow U_1^2 + U_2^2 \)
3. if \( X > 1 \) go to step (1)
4. Set \( Y \leftarrow \sqrt{-\frac{2\text{log}(X)}{X}} \)
5. return \( (Z_1 = U_1 Y, Z_2 = \leftarrow U_2 Y) \)

Step 3 ensures that why are not sampling outside the unit circle
and since \( -2\text{log}(X) \sim \exp(\frac{1}{2}) \) we know that \( \sqrt{-2\text{log}(X)} = \sqrt{R} \),
like in the Box-Muller method.
Finally \( (\frac{U_1}{\sqrt{X}}, \frac{U_2}{\sqrt{X}}) \) is uniformly distributed on the unit circle
and acts like \( \sin, \cos \).



Marsaglia-Bray-Algorithm 

Explain how affine transformations affect multivariate normal distributions

Let \( x ~ \mathcal{N}(\overline{x}, \mathbf{C_y}) \) be a normally distributed multivariate.
The new random variable \( \mathbf{y} = Ax + b \) obtained via an affine transformation is again normally distributed
with \( \mathbf{y} \sim \mathcal{N}(\mathbf{A}\mathbf{\bar{x}+\mathbf{b}, \mathbf{A}\mathbf{C_x}\mathbf{A}^\top} \).

In particular this reads:
\[
f_\mathbf{Y}(\mathbf{y)}
= {1 \over \sqrt{\lvert2\pi\mathbf{A}\mathbf{C_x}\mathbf{A}^\top\rvert}}
\exp\left(- {1 \over 2} \big[\mathbf{y}-(\mathbf{A}\mathbf{\bar{x}}+\mathbf{b}) \big]^\top (\mathbf{A}\mathbf{C_x}\mathbf{A}^\top)^{-1} \big[\mathbf{y}-(\mathbf{A}\mathbf{\bar{x}}+\mathbf{b}) \big] \right)
\]

The mean is induced by
\(
\mathbf{\bar{y}} = \mathbb{E}\{\mathbf{y}\} = \mathbb{E}\{\mathbf{A}\mathbf{x}+\mathbf{b}\} = \mathbf{A}\mathbb{E}\{\mathbf{x}\}+\mathbf{b} = \mathbf{A}\mathbf{\bar{x}}+\mathbf{b}
\)
and the covariance by
\[
\begin{array}{rcl}
\mathbf{C_y} & \triangleq & \mathbb{E}\{(\mathbf{y}-\mathbf{\bar{y}})(\mathbf{y}-\mathbf{\bar{y}})^\top\} \\
& = & \mathbb{E} \Big\{ \Big[ (\mathbf{A}\mathbf{x}+\mathbf{b})-(\mathbf{A}\mathbf{\bar{x}}+\mathbf{b}) \Big] \Big[ (\mathbf{A}\mathbf{x}+\mathbf{b})-(\mathbf{A}\mathbf{\bar{x}}+\mathbf{b}) \Big] ^\top \Big\} \\
& = & \mathbb{E} \Big\{ \Big[ \mathbf{A}(\mathbf{x}-\mathbf{\bar{x}}) \Big] \Big[ \mathbf{A}(\mathbf{x}-\mathbf{\bar{x}}) \Big] ^\top \Big\} \\
& = & \mathbb{E} \Big\{ \mathbf{A}(\mathbf{x}-\mathbf{\bar{x}}) (\mathbf{x}-\mathbf{\bar{x}})^\top \mathbf{A}^\top \Big\} \\
& = & \mathbf{A} \mathbb{E} \Big\{ (\mathbf{x}-\mathbf{\bar{x}}) (\mathbf{x}-\mathbf{\bar{x}})^\top \Big\} \mathbf{A}^\top \\
& = & \mathbf{A}\mathbf{C_x}\mathbf{A}^\top
\end{array}
\]


(This matters to the Cholesky decomposition as well as principal component analysis)

What is a Cholesky factorization, and how is it used for multivariate normals?

If we're able to sample multivariate \( Z \sim \mathcal{N}(\mathbf{0}, \mathbf{I}) \) then the Cholesky factorization
allows us an efficient generation of multivariate normals with arbitrary mean \( \mu \)and positive definite covariance \( \Sigma = AA^{T}\),
where \( A \) is a lower triangular matrix with real and positive diagonal entries.

By the affine transformation law for multivariate normals, we know that \( X = \mu + AZ \sim \mathcal{N}(\mu, AIA^{T}) \).

\( X \) is then given by 
\( X_1 = \mu_1 + A_{11} Z_1 \)
\( X_2 = \mu_2 + A_{21} Z_1 + A_{22} Z_{2}\)
\( X_d = \mu_d + A_{d1} Z_1 + A_{d2} Z_{2} \dots + A_{dd}Z_d\)

and \( A \) can actually be recursively be determined
\( \Sigma_{ij} = \sum{k+1}^{j} A_{ik}A_{jk}, j \leq i \), so \( A_{ij} = \frac{\Sigma_{ij}- \sum^{j-1}_{k=1}A_{ik}A_{jk}}{A_{jj}}, j < i \)
and \( A_{ii} = \sqrt{\Sigma_{ii} - \sum{k=1}^{i-1} A^2_{ik}} \)


Define stationarity and invariance for a time-homogeneous Markov process. How
are these notions related?

Actually in the lecture there was no requirement for time-homogeneity in the succeedings: 
For an initial distribution \( \nu \) let \( \nu_t = \nu \mu_t \) be the marginal distribution of \( X_t, t \in T \).
1. A probability measure \( \nu \) on \( (S, \mathcal{S}) \) is invariant, if \(\nu_t = \nu, \forall t \in T\)
2. The Markov process \( X \) is stationary, if \( \mathcal{L}(\vartheta_t X = \mathcal{L}(X), \forall t \in T \), i.e., \( \mathbb{P}^\nu \) is invariant under shifts
\( \veta_t : S^T \to S^T \)

The shift operator(s) acts on paths \( \omega : T \to S \) as \( \theta_t(\omega)(s) = \omega(t + s) \) is measurable and commutes
with the canonical construction, since its the identity.

Additionally we have the following proposition 
A Markov process is stationary with initial distribution \( \nu \), if and only if \( \nu \) is invariant for \( \mu_t \).

Define the following objects for time-homogeneous Markov process in discrete time:
hitting time, occupation time, associated probabilities, recurrence, transience.

Hitting time:

The initial hitting time is \( \tau_y = \inf_{n \in \mathbb{N}} X_n = y \)
and the \( k+1 \) hitting time \( \tau^{k+1}_y = \tau^{k}_y + \tau_y \circ \vartheta_{\tau^{k}_y} \)
by convention \( \tau_y^0 = 0 \).

Occupation time
\( \kappa_y = \text{argmax}_{k \in \mathbb{N}} \tau^k_y = \sum_{n\geq 1} \mathbb{1}_{X_n = y} \)
(so we are interested in \( k \) not in \( \tau^k \)).

The associated probabilities are

\( r_{xy} = \mathbb{P}_x[\tau_y < \infty] = \mathbb{P}_x[\kappy_y > 0] \)
in particular
\( \mathbb{P}_x[\kappa_y \geq \infty] = \mathbb{P}_x[\kappa^k_y < \infty] = r_{xy}r_{yy}^{k-1}\)
\( \mathbb{E}_x[\kappa_y] = \frac{r_{xy}}{1 - r_{yy}} \)

Show that a sufficient condition for recurrence is the positivity of an invariant
distribution.

TODO ??? Not in the lecture, the definitions also do not seem right

In the discrete time setting, transition kernels can be represented by (infinite)
matrices; explain how. How do the Chapman-Kolmogorov relations look like in
this notation?

Let \( S \) be countable and \( T = \mathbb{N}_{0} \), then 
the so called transition matrix is an \( |S|^2 \) matrix \( M \) 
(or homomorphism should \( |S| = \infty \))
with entries \( p_{i,j} = \mathbb{P}[X_{n+1} = j|X_n = i] \), \( i, j \in S \)
and the \( n \)-step probability is then \( p^{n}_{i,j} = M^{n}_{i,j} \).

The entries must satisfy
1. \( p_{i,j} \geq 0, \forall i, j \in S\)
2. \( \sum_{j \in S} p_{i,j} = 1\)


The Chapman-Kolmogorov relations are
\( p^{m+n}_{i,k} = \sum_j p_{i,j}^m p_{j,k}^n \)
or in matrix notation \( M^{m+n} = M^{m}M^n \)

Calculating \( e_i M \) (where \( e_i \) is the \( i \)-th unit vector) gives the vector  
\( (p^n_{i,1}, p^n_{i,2}, \dots, p^n_{i,|S|}) \) that is the vector with the entry in the  
\( j \)-th position as the probability of moving to state \( j \) in \( n \) steps from state \( i \).

A row vector \( \pi = (\pi_1, \pi_2, \dots, \pi_{|S|} \) is called an 
is called an equilibrium or stationary distribution of the Markov chain if
1. \( \pi_i \geq 0 \) for all states in \( S \)
2. \( \sum_{i \in S} = 1 \)
3. \( \pi P = \pi \)

otherwise we could interpret a vector \( \pi \) as the initial distribution \( \nu \)
of our \( X_0 \).

What is the period of a state? What is an irreducible Markov chain? State some
facts on recurrence, transience, period, invariant measures for irreducible Markov
chains.

The period \( d_i \) of state \( i \) is the greatest common divisor of \( \{n \in \mathbb{N} : p^n_{i,i} > 0\} \)
A state is aperiodic if \( d_i = 1 \)

If \( i \in S \) has period \( d \), then \( p_{i,i}^{nd} > 0\) for all but finitely many \( n \).

A Markov chain is irreducible, if
\( r_{i,j} = \mathbb{P}_i[\tau_j < \infty] > 0, \forall i, j \in S \)
that is at any point in time there is another, where the probability of moving into state 
\( j \) from state \( i \) is non-zero.

Such irreducible Markov chains satisfy
1. The states are either all recurrent or transient.
2. All states have the same period.
3. If \( \nu \) is invariant, then \( \nu_i > 0 \) for all \( i \in S \).

State the ergodic theorem for irreducible Markov chains.

Define the total variation of a signed measure of a space \( (S, \mathcal{S}) \)as \( \|\mu\| = \sup_{A \in \mathcal{S}} |\mu(A)| \) 
The ergodic theorem then states

For any irreducible, aperiodic Markov chain in \( S \) holds either (1) or (2):
1. There exists a unique invariant distribution \( \nu \) with \( \nu_i > 0 \) for all \( i \in S \), 
and for any initial distribution \( \mu \) on \( S \):
\( \lim_{n \to \infty} \|\mathbb{P}_\mu \circ \vartheta^{-1}_n - \mathbb{P}_\nu \| = 0\)
2. No invariant distribution exists, and we have
\( \lim_{n \to \infty} p^n_{i,j} = 0 \), \( i, j \in S \)


Recurrence and Transcience
\( \mathbb{P}_{x}[\kappa_x \geq k] = r_{x x}^k \)

Recurrence : \( r_{x x} = 1 \)
The numbers of return to x is almost surely infinite.

Transcience \( r_{x x} < 1\)
Since \( \sum_k \mathbb{P}x[\kappa_x \geq k] (1 - r_{x x})^{-1} - 1 \) we have \( \mathbb{P}_{x}[\kappa_x = \infty] = 0\)
by the Borel-Cantelli Lemma. One computes \( \kappa_x \) posseses a geometric distribution under \( \mathbb{P}_x \) i.e.
\( \mathbb{P}_x[\kappa_x = k] = r^k_{x x}(1 - r_{x x})\)
If the Markov process possesses an invariant distribution \( \nu \) with \( \nu(x) > 0 \), then \( x \) is recurrent.
    
